\section{最坏情况下的最优与交换次序}
\label{sec:17-Constrained-Guarantees/06-Games-and-Equilibria/01}

一家支付机构有两条入金通道 A 和 B。风控每天只有人力盯住其中一条，欺诈团伙每天也只走其中一条。把一天的净结果折成分数记在防守方账上：查 A 且对方走 A，拦下来记 \(+3\)；查 A 而对方走 B，白查一趟又放过一笔，记 \(-1\)；查 B 而对方走 A，A 通道的单笔额度更大，记 \(-2\)；查 B 且对方走 B，B 通道的欺诈金额更高，拦下来记 \(+4\)。

\[
A=\begin{pmatrix}3 & -1\\ -2 & 4\end{pmatrix},
\]

行是防守方的两个选择，列是攻击方的两个选择，数值是防守方的收益、同时是攻击方的损失。

风控主管问的问题很实际：不管对方怎么走，我最少能保证拿到多少。查 A 的最坏情况是 \(-1\)，查 B 的最坏情况是 \(-2\)，所以保底值是 \(-1\)，办法是固定查 A。写成式子就是 \(\max_x\min_y u(x,y)=-1\)。

攻击方那边的主管问同一个问题，只是符号相反：不管风控怎么查，我最多会被扣多少。走 A 时防守方最好的应对给出 \(3\)，走 B 时给出 \(4\)，所以他的保底是 \(3\)，办法是固定走 A。写成式子是 \(\min_y\max_x u(x,y)=3\)。

两个数不一样，而且差得不小：\(-1\) 和 \(3\)。同一个矩阵，同一个问法，答案取决于谁先说话。

\subsection{不等号的方向是固定的}

这不是这个矩阵的巧合。对任意的 \(u\)，只要两边的极值都存在，就有

\[
\max_x\min_y u(x,y)\;\le\;\min_y\max_x u(x,y).
\]

证明一行。对任意一对 \((x,y)\)，把 \(u(x,y)\) 上下夹住：

\[
\min_{y'}u(x,y')\;\le\;u(x,y)\;\le\;\max_{x'}u(x',y).
\]

左端只与 \(x\) 有关，右端只与 \(y\) 有关，而这条链对所有 \((x,y)\) 成立。于是先对 \(x\) 取最大——左端变成 \(\max_x\min_{y'}u(x,y')\)，右端不受影响；再对 \(y\) 取最小——右端变成 \(\min_y\max_{x'}u(x',y)\)，左端不受影响。两头合上即得。

不等式的含义是后动的一方占便宜。左边那个数是这样得来的：我先把策略摆在桌面上，对方看清之后再挑最狠的应对，我在这种条件下能拿到的最好结果。右边那个数是对方先摆、我后挑。谁先亮牌谁吃亏，差额 \(3-(-1)=4\) 就是先手信息的价值。

这也说明了为什么``最优策略''在博弈里不是一个自明的概念。单人优化里 \(\max_x f(x)\) 只有一个含义，因为 \(f\) 是死的；有对手时目标函数带着对方的变量，先取哪个极值就成了必须交代清楚的事。

\subsection{允许随机化之后两个数合一}

那条不等式一般是严格的。使它变成等式的，是 von Neumann 的极小极大定理：把纯策略换成混合策略之后，两侧相等。

混合策略就是在自己的选项上放一个概率分布。防守方的策略集从两个点变成单纯形 \(\Delta_2=\{p\in\R^2:p_i\ge0,\;p_1+p_2=1\}\)，攻击方同理，效用变成双线性的 \(u(p,q)=p^{\trans}Aq\)。

在这个例子里两侧的值都能手算。设防守方以概率 \(p\) 查 A。对方走 A 时的期望收益是 \(3p-2(1-p)=5p-2\)，走 B 时是 \(-p+4(1-p)=4-5p\)。防守方要让两者中的较小者最大，而一升一降，最优点在两者相等处：\(5p-2=4-5p\) 给出 \(p^\star=0.6\)，代回得 \(1\)。设攻击方以概率 \(q\) 走 A，防守方查 A 的期望是 \(3q-(1-q)=4q-1\)，查 B 是 \(-2q+4(1-q)=4-6q\)，攻击方要让较大者最小，得 \(q^\star=0.5\)，代回同样得 \(1\)。

\[
\max_{p}\min_{q}p^{\trans}Aq=\min_{q}\max_{p}p^{\trans}Aq=1 .
\]

原来的两个数 \(-1\) 与 \(3\) 之间，出现了一个唯一的 \(1\)。这个数叫这个博弈的值：防守方能保证不低于它，攻击方能保证不高于它，双方都无法再改善。先手的优势消失了——因为混合策略把``我的选择''藏了起来，对方即使知道 \(p^\star=0.6\) 这个分布，也不知道今天这一次会落在哪一边。

\begin{center}
\begin{tikzpicture}[x=1cm,y=1cm,font=\small,>=stealth]
  \draw[black!55] (1.0,1.6) rectangle (3.6,3.4);
  \draw[black!55] (2.3,1.6) -- (2.3,3.4);
  \draw[black!55] (1.0,2.5) -- (3.6,2.5);
  \node at (1.65,2.95) {\(3\)};
  \node at (2.95,2.95) {\(-1\)};
  \node at (1.65,2.05) {\(-2\)};
  \node at (2.95,2.05) {\(4\)};
  \node[font=\footnotesize,black!70] at (1.65,3.66) {攻 A};
  \node[font=\footnotesize,black!70] at (2.95,3.66) {攻 B};
  \node[font=\footnotesize,black!70] at (0.5,2.95) {查 A};
  \node[font=\footnotesize,black!70] at (0.5,2.05) {查 B};
  \node[font=\footnotesize,blue!65!black] at (4.25,3.66) {行最小};
  \node[font=\footnotesize,blue!65!black] at (4.25,2.95) {\(-1\)};
  \node[font=\footnotesize,blue!65!black] at (4.25,2.05) {\(-2\)};
  \node[font=\footnotesize,red!70!black] at (0.5,1.24) {列最大};
  \node[font=\footnotesize,red!70!black] at (1.65,1.24) {\(3\)};
  \node[font=\footnotesize,red!70!black] at (2.95,1.24) {\(4\)};
  \node[anchor=west,align=left,font=\footnotesize] at (0.35,0.45)
    {\(\max_x\min_y=-1\)，\(\min_y\max_x=3\)};

  \begin{scope}[shift={(6.2,2.7)}]
    \draw[->,black!60] (0,0) -- (2.9,0) node[right,font=\footnotesize] {\(p\)};
    \draw[->,black!60] (0,0) -- (0,2.2);
    \draw[black!35,dashed] (0,0.99) -- (2.9,0.99);
    \draw[black!35,dashed] (1.56,0) -- (1.56,0.99);
    \draw[blue!65!black,thick] (0,0) -- (1.56,0.99) -- (2.6,0.33);
    \fill[blue!65!black] (1.56,0.99) circle (1.6pt);
    \node[anchor=west,font=\footnotesize,blue!65!black] at (1.05,1.78)
      {\(\min_q p^{\trans}Aq\)};
    \node[anchor=north,font=\scriptsize,black!70] at (1.56,-0.06) {\(p^\star=0.6\)};
    \node[anchor=east,font=\scriptsize,black!70] at (-0.05,0.99) {\(1\)};
  \end{scope}

  \begin{scope}[shift={(6.2,-0.3)}]
    \draw[->,black!60] (0,0) -- (2.9,0) node[right,font=\footnotesize] {\(q\)};
    \draw[->,black!60] (0,0) -- (0,2.2);
    \draw[black!35,dashed] (0,0.99) -- (2.9,0.99);
    \draw[black!35,dashed] (1.3,0) -- (1.3,0.99);
    \draw[red!70!black,thick] (0,1.98) -- (1.3,0.99) -- (2.6,1.65);
    \fill[red!70!black] (1.3,0.99) circle (1.6pt);
    \node[anchor=west,font=\footnotesize,red!70!black] at (1.35,1.92)
      {\(\max_p p^{\trans}Aq\)};
    \node[anchor=north,font=\scriptsize,black!70] at (1.3,-0.06) {\(q^\star=0.5\)};
    \node[anchor=east,font=\scriptsize,black!70] at (-0.05,0.99) {\(1\)};
  \end{scope}
\end{tikzpicture}

\emph{左边是纯策略下的两个不同答案 \(-1\) 与 \(3\)；右边两条折线的极值点落在同一条虚线上，允许随机化之后这两个数合并成了一个 \(1\)。}
\end{center}

图里右上那条折线是防守方的保底收益随 \(p\) 的变化——它是若干条直线取下确界，所以是凹的折线；右下那条是攻击方视角的同一件事，若干条直线取上确界，所以是凸的折线。凹的那条的最高点与凸的那条的最低点重合，这就是极小极大定理在两维情形下的全部内容。折线的凹凸不是画得巧，是``一族线性函数取下确界必凹''这条构造规则的直接结果，见\hyperref[sec:09-Convex-Optimization/03-Convex-Functions/04]{``怎样识别和构造凸函数''}。

\subsection{这是``两个操作能不能交换''的又一个实例}

极小极大定理表面上讲博弈，骨子里讲的是次序：什么时候 \(\max\) 与 \(\min\) 可以对调。\hyperref[sec:18-Synthesis/02-Recurring-Ideas/05]{``换序：两个操作能不能交换顺序''}把这个动作单独拎出来讲过，本书别处的几处出现值得并排看一眼。

\hyperref[sec:12-Real-Analysis-Support/04-Limits-and-Integration/01]{``点态极限为何不能直接穿过积分号''}里的反例是 \(f_n=n\ind_{(0,1/n]}\)：逐点极限处处为零，每个 \(f_n\) 的积分却恒为 \(1\)。质量没有消失，它在取极限的过程中跑到了一个越来越窄、越来越高的地方，极限函数看不见它。\hyperref[sec:12-Real-Analysis-Support/04-Limits-and-Integration/03]{``求导穿过积分号需要控制差商''}防的是同一件事，只是逃走的是差商。

这里逃走的是最优反应。先固定 \(x\) 再让对方挑，对方的挑法依赖 \(x\)；先固定 \(y\) 再让我挑，我的挑法依赖 \(y\)。两条路径上被取极值的对象不是同一个，所以结果不必相同。让它们相同需要一个额外的结构——某个鞍点必须真的存在，也就是存在一对 \((p^\star,q^\star)\)，使得任一方单独偏离都不会变好。上面那对 \((0.6,0.5)\) 就是。

充分条件是策略集凸紧、效用连续且对自己的变量凹、对对方的变量凸。有限矩阵博弈自动满足全部四条：单纯形是凸紧集，双线性函数对两个变量分别既凹又凸。Sion 把条件放松到拟凹与拟凸、并且只需一侧紧，但结构是一样的——凸性负责让两条路径的极值点碰头，紧性负责让极值真的取得到。

\hyperref[sec:09-Convex-Optimization/05-Duality/02]{``对偶间隙何时闭合''}讲的是同一件事换了衣服。Lagrangian 的对偶间隙就是 \(\inf_x\sup_\lambda L\) 与 \(\sup_\lambda\inf_x L\) 之差，强对偶成立等于这两个操作可以交换。矩阵博弈的极小极大定理与线性规划的强对偶定理可以互相推出，它们是同一个结论的两种叙述。\hyperref[sec:09-Convex-Optimization/05-Duality/01]{``Lagrangian 怎样制造全局下界''}里那句``弱对偶永远成立，强对偶需要条件''，与这里的``不等式永远成立，等式需要条件''是逐字对应的。

条件塌掉的样子也值得看。策略集不紧时定理会失效：允许攻击方从开区间 \((0,1)\) 里选一个数，最优反应可能在端点处取不到，两侧的值就分开了。凹凸性去掉之后同样失效，一般的连续函数没有理由拥有鞍点。

还有一条不写在数学条件里、却在实践中最先塌掉的：随机化必须是不可预测的。定理里的对手只知道分布 \(p^\star\)，不知道这一次抽到了哪个。若对方能观察到实现值——伪随机种子泄露、抽签流程可被内部人员看到、甚至只是执行者有可辨认的习惯——混合策略就退化回纯策略，博弈的值从 \(1\) 掉回 \(-1\)。这一条与\hyperref[sec:17-Constrained-Guarantees/04-Online-Learning-and-Regret/03]{``探索与利用是同一笔预算''}里那条``强制探索必须真的随机''是同一个要求。

\subsection{为什么必须允许随机}

石头剪刀布是最短的说明。三个纯策略，任何一个都有唯一的克制者，最优反应绕着圈跑：对方出石头我该出布，我出布对方该出剪刀，对方出剪刀我该出石头，回到起点。纯策略集合里不存在任何一对策略能互为最优反应，鞍点根本不存在。

把随机化允许进来，策略集从三个孤立的点变成一个二维单纯形——一个凸紧集，见\hyperref[sec:09-Convex-Optimization/02-Convex-Sets/01]{``线段、凸组合与凸包''}。均匀分布 \((1/3,1/3,1/3)\) 是鞍点，博弈的值是 \(0\)。

关键在于这不是一个取巧的补丁。定理要的是凸紧的策略集与线性的效用，而有限选项的集合两条都不满足：三个点不是凸集，在它上面定义的函数也谈不上凹凸。混合策略做的事情正是把这个集合的凸包取出来当作新的可行集，同时把效用延拓成期望——期望关于概率向量是线性的，凹凸性于是自动成立。\hyperref[sec:09-Convex-Optimization/03-Convex-Functions/01]{``弦在图像上方意味着什么''}里那句``凸性保证从一个可行点走向另一个可行点的路上不会出意外''，在这里读作``两个策略可以按比例掺着用''。

随机化不是技巧，是为了让可行集变凸。这句话在本书里出现过不止一次：凸松弛把整数变量放到区间上，也是为了同一件事。

\subsection{零和是特例}

上面所有结论都依赖一个假设：\(u\) 是一个数，一方的收益就是另一方的损失。有了这个假设，才谈得上``博弈的值''这样一个双方共享的量，才谈得上 \(\max\min\) 与 \(\min\max\) 是同一个函数的两种取法。

真实的博弈很少这么干净。两家公司同时降价，双方都变差；两条道路同时拥堵，所有司机都变差；两个团队互相拆台，公司整体受损。这些情形里每个参与者有各自的收益函数，不存在一个可以对调 \(\max\) 与 \(\min\) 的单一目标，极小极大定理无从谈起。

一旦丢掉零和，还剩下什么可说？剩下的是一个更弱、也更普遍的概念：一组策略，使得每个人在别人不动的前提下都无法单方面改善。它不再是某个数的最优值，而是一个平衡条件。下一节说明这样的平衡为什么一定存在，以及它的存在性靠的是本书里另一个反复出现的动作。
